\section{校优摘要的组织原则}

校级优秀毕业实习与设计（论文）摘要不是主论文摘要页的简单复制，而是一篇围绕研究价值、方法选择、结果证据和结论贡献重新组织的独立短文。正文通常应在有限页数中交代研究背景、关键问题、核心方案和验证结果，并避免把原论文目录式地压缩粘贴。为此，模板将首屏信息与正文结构拆开：题名、作者信息和中文摘要位于首页顶部，后续内容直接进入双栏正文，以便在固定版心中呈现更高的信息密度。

围绕摘要写作目标，研究者应优先保留可以支撑评价判断的核心信息，例如问题场景、方案差异、实验或设计结果、结论解释以及局限说明。相关工作和理论背景可以适度压缩，但不应完全省略，否则摘要将难以说明课题定位。已有研究表明，面向评审的摘要文本如果同时兼顾问题定义、方法路径和结果证据，通常更有利于形成稳定判断\cite{example2026,examplebook}。

在版面安排上，校优摘要的第一页往往同时承担“快速说明课题价值”和“建立阅读节奏”两项任务，因此首节不宜过短。若一开篇只给出极少量背景信息，页面右栏就会出现较大的留白，视觉上容易显得结构松散，也不利于和后续图表、二级标题形成稳定的阅读密度。更稳妥的做法，是先用一到两段文字交代研究对象、研究边界与评价口径，再进入更细的分级标题。

\subsection{内容压缩与证据保留}

压缩并不意味着删减到只剩结论。校优摘要的主要难点在于：一方面需要保留与主论文同等层级的关键信息，另一方面又必须在有限篇幅内保持逻辑闭合。因此，摘要应优先保留“问题是什么、为什么重要、做了什么、得到什么、说明什么”这五类内容，并通过图表或数据提示增强可读性。对不直接影响结论判断的细节，可在正文中采用概述性语言处理，而不是完全展开。

如果正文第一页还需要放入图示或流程框图，则更应提前为右栏准备足够的文字体量。否则页面会呈现“左栏信息集中、右栏启动偏晚”的观感，虽然从排版机制上并非错误，但与官方示例页追求的紧凑双栏视觉并不完全一致。因此，示例模板在这里故意增加一段过渡性说明，用来模拟真实申报稿中更常见的版面密度。

\subsubsection{研究目标}
摘要中的研究目标应尽量具体，避免只写“提高效率”或“优化设计”这一类泛化表述。更稳妥的写法是说明对象、约束和评价指标，使评审能够迅速判断课题边界。

\subsubsection{结果表达}
结果表述应对应研究目标，尽量给出可比较、可复核的定量证据。如图\ref{fig:excellent-placeholder}所示，图示内容未必复杂，但应能够服务于论证，而不是仅作为装饰。对无法展开说明的复杂数据，可通过结论性归纳突出结果意义\cite{exampleconf}。

\begin{center}
    \fbox{\rule{0pt}{2.4cm}\rule{0.82\columnwidth}{0pt}}
    \captionof{figure}{校优摘要图示占位示例}
    \label{fig:excellent-placeholder}
\end{center}

\section{版式复用与独立成稿}

本模板采用“共享元数据 + 独立正文内容”的组织方式。共享元数据用于维护题名、作者、指导教师、学院专业等稳定字段，使主论文和校优摘要稿在基础信息上保持一致；而摘要正文、关键词和英文摘要则单独维护，以避免因全文结构差异而互相牵制。这样的拆分既满足复用需求，也符合校优摘要需要独立提炼的官方意图。

\subsection{双栏正文的使用边界}

双栏版式能够在有限页数内承载更多信息，但也会放大排版噪声。正文段落宜保持紧凑，句子长度不宜过度膨胀，图表应尽量控制在单栏内；若表格较宽，可考虑改写为简化表或在最终稿中使用跨双栏布局。表\ref{tab:excellent-outline}给出了一个适合校优摘要的内容组织示例。

\begin{center}
    \captionof{table}{校优摘要内容组织示例}
    \label{tab:excellent-outline}
    \begin{tabular}{p{0.18\columnwidth}p{0.72\columnwidth}}
        \toprule
        模块 & 建议保留的信息 \\
        \midrule
        中文摘要 & 问题、方法、结果、结论四类核心信息 \\
        正文第1节 & 研究背景、目标与课题价值 \\
        正文第2节 & 技术路线、实现过程或分析方法 \\
        正文第3节 & 关键结果、结论解释与应用意义 \\
        \bottomrule
    \end{tabular}
\end{center}

\subsubsection{共享字段}
适合共享的信息通常是不会因为摘要重写而改变的事实性字段，例如姓名、学院、专业、导师、英文题名和毕业届别。这类信息统一维护后，可以明显降低重复修改带来的遗漏风险。

\subsubsection{独立字段}
中文摘要、英文摘要、关键词、图表与作者简介更适合作为校优摘要稿的独立内容，因为它们直接服务于本稿件的评审目标，往往需要相较主论文版本做再组织和再裁剪。

\section{提交流程建议}

建议在主论文基本定稿后再启动校优摘要稿写作。此时研究问题、主要结果和参考文献范围已经稳定，更适合集中提炼。写作顺序上，可先列出五到七条最想让评审记住的关键信息，再据此生成中文摘要和正文骨架，最后补齐英文摘要、作者简介和图表说明。这样可以避免先写长文再被动压缩，提升最终稿的一致性和可读性。

\subsection{最终检查清单}

\subsubsection{内容完整性}
检查题名、摘要、关键词、正文、参考文献、作者简介和英文摘要是否全部出现，确保结构完整且信息闭环。

\subsubsection{提交友好性}
检查首页标题区、双栏正文、图表位置和页眉页脚是否稳定；同时核对关键词分隔、参考文献数量和作者简介内容，确保正式提交前不留说明性占位文本。
